\subsection{CHARACTERISTIC POLYNOMIAL OF AN ENDOMORPHISM}

Let M be a free A-module of dimension n and u an endomorphism of M. Consider the polynomial ring in two indeterminates A{[}X, Y{]} and the A{[}X, Y{]}-module M{[}X, Y{]} = A{[}X, Y{]} ⊗\_A M; let \(\bar{u}\) be the endomorphism of the A{[}X, Y{]}-module M{[}X, Y{]} canonically derived from u (II, §5, no. 1). It follows from no. 5, Proposition 11 that

\[
\det (\mathbf {X} - \mathbf {Y} \bar {\boldsymbol {u}}) = \sum_ {j = 0} ^ {n} (- 1) ^ {j} \operatorname{Tr} \left(\bigwedge^ {j} (u)\right) \mathbf {X} ^ {n - j} \mathbf {Y} ^ {j}\tag{49}
\]

for if \(U\) is the matrix of \(u\) with respect to a basis \((e_i)_{1 \leqslant i \leqslant n}\) of \(\mathbf{M}\) , \(U\) is the matrix of \(\bar{u}\) with respect to the basis \((1 \otimes e_i)_{1 \leqslant i \leqslant n}\) of \(\mathbf{M}[\mathbf{X}, \mathbf{Y}]\) , hence

\[
\operatorname{Tr} \left(\bigwedge^ {j} (\bar {u})\right) = \operatorname{Tr} \left(\bigwedge^ {j} (u)\right).
\]

DEFINITION 3. Let M be a finite-dimensional free A-module and u an endomorphism of M. The determinant of the endomorphism X - \(\bar{u}\) of the free A{[}X{]}-module M{[}X{]} is called the characteristic polynomial of u and is denoted by \(\chi_u(X)\) .

If \(\mathbf{M}\) is of rank \(n\) , it follows from (49) that

\[
\chi_ {u} (\mathbf {X}) = \sum_ {j = 0} ^ {n} (- 1) ^ {j} \operatorname{Tr} \left(\bigwedge^ {j} (u)\right) \mathbf {X} ^ {n - j}\tag{50}
\]

for \(\operatorname{det}(\mathbf{X} - \mathbf{Y}\bar{u}) = \operatorname{det}(\mathbf{X}.I_n + \mathbf{Y}U)\) and \(\operatorname{det}(\mathbf{X} - \bar{u}) = \operatorname{det}(\mathbf{X}.I_n - U)\) . It is therefore seen that \(\chi_u(\mathbf{X})\) is a monic polynomial of degree \(n\) , in which the coefficient of \(\mathbf{X}^{n-1}\) is \(-\mathrm{Tr}(u)\) and the constant term is \((-1)^n \operatorname{det}(u)\) .

PROPOSITION 20 (``Cayley-Hamilton theorem''). For every endomorphism \(u\) of a finite-dimensional free A-module, \(\chi_u(u) = 0\) .

In the notation of Proposition 18 (§ 3, no. 10), for all \(x \in \mathbf{M}\) , \(\chi_u(u)(x)\) is the image under \(\phi\) of \(\chi_u(\mathbf{X}) \otimes x\) . But if \(v\) is the endomorphism of \(\mathbf{M}[\mathbf{X}]\) , the co-transpose of \(\mathbf{X} - \bar{u}\) (no. 6), then

\[
\chi_ {u} (\mathbf {X}) \otimes x = \chi_ {u} (\mathbf {X}) (1 \otimes x) = (\mathbf {X} - \bar {u}) (v (1 \otimes x))
\]

and the conclusion follows from Proposition 18 of no. 10.

